\documentclass[UTF8]{ctexart}
\usepackage{xeCJK}
\usepackage{geometry}
\geometry{left=2cm,right=2cm,top=2.5cm,bottom=2.5cm}
\setlength{\headheight}{12.7pt}

\usepackage{titlesec}
\titleformat{\section}
  {\normalfont\large\bfseries}
  {\thesection}
  {1em}
  {}
\titlespacing*{\section}
  {0pt}
  {3.5ex plus 1ex minus .2ex}
  {2.3ex plus .2ex}

\usepackage{amsmath}
\usepackage{amssymb}
\usepackage{amsthm}
\usepackage{enumitem}
\setlist[enumerate]{itemsep=0pt,topsep=0pt,parsep=0pt,leftmargin=0pt,itemindent=2.5em}
\setlist[itemize]{itemsep=0pt,topsep=0pt,parsep=0pt,leftmargin=0pt,itemindent=2.5em}
\usepackage{fancyhdr}
\pagestyle{fancy}
\fancyhf{}
\usepackage{graphicx}
\usepackage{hyperref}
\usepackage{indentfirst}
\usepackage{lmodern}


\begin{document}
\setlength{\abovedisplayskip}{1pt}
\setlength{\belowdisplayskip}{1pt}

\section{网}

\hypertarget{sec:定义1.1}{}
\noindent\textbf{定义1.1 (网)}

给定\href{../../01 基础集合论/04 序理论/01 序关系/01 预序关系#nameddest=sec:定义2.1}{有向集}
$(D,\leqslant)$和集合$X$, 称映射$\varphi:D\to X$为$X$中的一个\textbf{网}.

\vspace{10pt}

\hypertarget{sec:定义1.2}{}
\noindent\textbf{定义1.2 (网的尾部)}

给定网$\varphi:D\to X$和某个$\alpha_0\in D$, 定义
\[
    T_{\alpha_0}=\{\varphi(\alpha)\mid\alpha\geqslant\alpha_0\},
\]
称之为$\varphi$在$\alpha_0$之后的\textbf{尾部}.

\vspace{10pt}

\hypertarget{sec:定义1.3}{}
\noindent\textbf{定义1.3 (子网)}

给定网$\varphi:D\to X$和$\psi:E\to X$. 如果存在映射$f:E\to D$使得:
\begin{enumerate}
    \item $f$保序, 即$\forall\beta,\beta'\in E:\beta\leqslant\beta'\to
    f(\beta)\leqslant f(\beta')$,
    \item $f$是\href{../../01 基础集合论/04 序理论/01 序关系/01 预序关系#nameddest=sec:定义3.2}
    {共尾映射}, 即$\forall\alpha\in D\,\,\exists\beta\in E:f(\beta)\geqslant\alpha$,
    \item $\psi=\varphi\circ f$,
\end{enumerate}
就称$\psi$是$\varphi$的\textbf{子网}.

\vspace{10pt}

给定$X$中的网$\varphi$和$U\subset X$. 如果$\varphi$的某个尾部包含于集合$U$, 即
\[
    \exists\alpha_0\in D\,\,\forall\alpha\geqslant\alpha_0:\varphi(\alpha)\in U,
\]
就说该网\textbf{最终落入}$U$; 如果$U$的原像在$D$中共尾, 即
\[
    \forall\alpha_0\in D\,\,\exists\alpha\geqslant\alpha_0:\varphi(\alpha)\in U,
\]
就说该网\textbf{频繁落入}$U$. 显然:
\begin{enumerate}
    \item 如果$\varphi$最终落入$U$, 那么它也频繁落入$U$,
    \item $\varphi$频繁落入$U$当且仅当它不最终落入$U^c$.
\end{enumerate}

\vspace{10pt}

\hypertarget{sec:定义1.4}{}
\noindent\textbf{定义1.4 (泛网)}

给定$X$中的网$\varphi$. 如果对任意的$U\subset X$,
$\varphi$都最终落入$U$或$U^c$, 就称它是\textbf{泛网}.






% \newpage

% \section{网的极限与聚点}

% \hypertarget{sec:定义2.1}{}
% \noindent\textbf{定义2.1 (网收敛)}

% 给定\href{../01 拓扑空间/01 拓扑空间#nameddest=sec:定义1.1}{拓扑空间}
% $(X,\mathcal{T})$, $X$上的网$(x_\alpha)_{\alpha\in\Lambda}$和点$x\in X$. 如果
% \[
%     \forall N\in\mathcal{N}_x\,\,\exists\alpha\in\Lambda:T_\alpha\subset N,
% \]
% 即$x$的任意邻域包含一个尾部, 就称$(x_\alpha)_{\alpha\in\Lambda}$
% \textbf{收敛于}$x$, 也称$x$是$(x_\alpha)_{\alpha\in\Lambda}$的一个\textbf{极限},
% 记作$x_\alpha\to x$.

\end{document}